\scanpage{90}
% archive sheet number: 89
% block text | lines 1-5
i.e. $\Gamma$ opère \uline{fidèlement} sur le corps des restes. On a déjà vu (comme conséquence du th. d'indépendance algébrique des automorphismes, \unsure{resp.} le lemme de Lang) que $H^1(\Gamma,A/mA^*)=\unsure{\{e)|(e)}$. Prouvons alors par récurrence sur $n\geq1$ que
% block display | line 6
\[
H^1(\Gamma,A/m^n A^*)=\{e\}
\]
% block text | lines 7-8
Pour ceci, on \deleted{suppose} le suppose fait pour $n\geq1$, on écrit la suite exacte
% block diagram | line 9
\diagram{\begin{tikzcd}[ampersand replacement=\&,column sep=small]
e\arrow[r,dash]\&N_n\arrow[r,dash]\&A/m^{n+1}A^*\arrow[r,dash]\&A/m^n A^*\arrow[r,dash]\&e
\end{tikzcd}}
% block text | lines 10-14
et on constate qu'en tant que $\Gamma$-groupe, $N_n$ est isomorphe à $m^{n+1}A/m^n A$, donc à un espace vectoriel sur $k=A/m$. \unsure{Donc|Or} il est (en tant que $\Gamma$-groupe) isomorphe à $k^N$, donc $H^1(\Gamma,k^N)=H^1(\Gamma,k)^N=0$ en vertu du \uline{lemme} $H^i(\Gamma,k)=0$ pour $i>0$.

% block text | line 15
La suite exacte
% block diagram | line 16
\diagram{\begin{tikzcd}[ampersand replacement=\&]
H^1(\Gamma,N_n)\arrow[r,dash]\&H^1(\Gamma,A/m^{n+1}A^*)\arrow[r,dash]\&H^1(\Gamma,A/m^n A)^*
\end{tikzcd}}
% block text | line 17
montre donc que l'on a
% block display | line 18
\[
H^1(\Gamma,A/m^{n+1}A^*)^{\deleted{\unsure{*}}}=\{e\}.
\]
